\section{An ordinary Eisenstein congruence over the rationals}

The rank-two construction needs a cuspidal representation in every
sufficiently large weight of one parity. It needs the order of the two
local characters, even when the residual Galois representation is
reducible. We construct this particular family over $\mathbb Q$.
All reductions are at the place specified by
$\iota:\overline{\mathbb Q}_\ell\simeq\mathbb C$.

% BEGIN P66E1
\paragraph{P66E1. The primitive tame character and its auxiliary prime.}
Use the data of $\mathrm{PropEGP}$. Thus $\ell$ is odd, $w$ is a
prime prime to $N$, $\ell\nmid N$, and the nontrivial $\theta$,
of order $\ell$, factors through the action on $\mu_w$.
The nontrivial $\phi$ factors through the action on $\mu_N$ and is
one on elements acting there by the $w$th power. Cyclotomic
reciprocity therefore writes $\phi,\theta$ as the Galois characters
of primitive Dirichlet characters $\phi_0,\theta_0$ of conductors
$M\mid N,w$, respectively. Indeed the only smaller divisor of the
prime $w$ is one, and $\theta$ is nontrivial. Similarly $M>1$.
Surjectivity of the action on roots of unity gives
$\phi_0(w)=1$, and $\theta_0$ has exact order $\ell$. Hence
$w\equiv1\pmod\ell$, $w\ne\ell$, and $\theta_0(-1)=1$.
The product $\eta=\phi_0\theta_0$ is primitive of conductor
$C=Mw$ by the Chinese remainder theorem. For a specified complex
conjugation $c$, the parity test $\phi(c)=(-1)^k$ is precisely
$\eta(-1)=(-1)^k$. The level $C$ is prime to $\ell$.
% END P66E1

% BEGIN P66E2
\paragraph{P66E2. Why the Eisenstein constant vanishes on reduction.}
Write $L_k(x)=(x\,d/dx)^{k-1}(x/(1-x))$ for $x\ne1$ and
$\zeta$ for a primitive $C$th root. For $k\ge3$ of the specified
parity, the normalized Eisenstein series $E=E_k(1,\eta)$ has
positive coefficient $a_n=\sum_{d\mid n}\eta(d)d^{k-1}$.
Its constant $b$ at infinity satisfies
\[
  2\,g(\eta^{-1})b
  =\sum_{a\bmod C}\eta^{-1}(a)L_k(\zeta^a).
\]
Terms with nonunits are omitted. In this formula the inverse Dirichlet
character is extended by zero. To see the normalization, form the
weight-$k$ Eisenstein sums for rows $(u,v)\equiv(0,a)\pmod C$,
weight by $\eta^{-1}(a)$ and divide the normalized sums by
$2g(\eta^{-1})$. Finite Fourier transformation gives the divisor
sum above: the contributions of the two signs agree by parity.
The zero-frequency term is the displayed rational function, by the
partial fraction expansion of the cotangent and its $k-1$
derivatives. After slashing by a matrix with lower left entry $c'$,
a row contributing to the constant must satisfy $a c'=0\bmod C$.
Thus constants vanish off the cusps above infinity of $X_0(C)$.

Each $L_k(\zeta^a)$ occurring with $\phi_0(a)\ne0$ is integral,
since its numerator is an integer polynomial in $\zeta^a$ and its
denominator $(1-\zeta^a)^k$ is a unit for $\zeta^a\ne1$ of
order prime to $\ell$. Put
$T=\sum_{b'\bmod M}\phi_0^{-1}(b')L_k(\zeta^{wb'})$,
again summing only units of the character. On $w\nmid a$, the
factor $\theta_0^{-1}(a)$ reduces to one. Summing over all $a$
with the $\phi_0^{-1}$ factor, including the multiples of $w$,
gives $w^kT$. Indeed grouping $a=b'+Mj$ uses
\[
 \sum_{j=0}^{w-1}L_k(x\zeta^{Mj})=w^k L_k(x^w).
\]
For $|x|<1$ this follows by summing $\sum_{n\ge1}n^{k-1}x^n$
over the roots of unity, and then by rational identity for $x^w\ne1$.
The multiples of $w$ alone contribute $T$, since $\phi_0(w)=1$.
Consequently $2g(\eta^{-1})b\equiv(w^k-1)T\equiv0$.
Both $2$ and the Gauss sum are units: the product with its inverse-character
Gauss sum is $\eta(-1)C$, a unit. This proves $b\equiv0$.
% END P66E2

% BEGIN P66E3
\paragraph{P66E3. A cusp eigenform in every sufficiently large admissible weight.}
An auxiliary modular form $U$ at the same weight, level and character
can be chosen integral, with constant one at infinity and zero
constants away from its diamond orbit. Here is a geometric way to
obtain the auxiliary form used in the source. Work on the tame
compactified modular curve $X_1(C)$ over a finite extension of
$\mathbb Z_\ell$ containing the cusp coordinates and character
values. Here $C=Mw\ge5$. The cusps form a disjoint finite
\'etale divisor after this extension. For large $k$, restriction
of $\omega^k$ to that divisor is surjective, by relative ampleness
of the Hodge bundle and Serre vanishing applied to
$\omega^k(-\mathrm{cusps})$. Lift the value one in the Tate
trivialization at the chosen infinity cusp and zero at all the
other cusps. Sum the diamond translates against the inverse
character. At the chosen cusp the only contributing diamonds are
$1,-1$, whose actions on the differential differ by $(-1)^k$.
Parity makes their contributions equal. Divide by $2$, an
$\ell$-adic unit. The resulting $U$ has the desired character
and constants. The integral $q$-expansion principle gives integral
coefficients. This argument uses the standard tame integral modular
curve and the ampleness of $\omega$ (Katz--Mazur,
\emph{Arithmetic Moduli of Elliptic Curves}, and
Deligne--Rapoport, \emph{Les sch\'emas de modules de courbes
elliptiques}). Equivalently the source constructs the diamond
translates using explicit modular forms.

Now $f=E-bU$ vanishes at all cusps. Indeed at every cusp in the
diamond orbit the two constants transform by the same character.
All positive coefficients of $f$ reduce to $a_n$ from P66E2 and
$a_1=1$. The divisor recurrence gives, for every prime $p$,
\[
 a_{pn}+\eta(p)p^{k-1}a_{n/p}
       =(1+\eta(p)p^{k-1})a_n ,
 \qquad a_{n/p}=0\quad(p\nmid n).
\]
For $p\mid C$ this is the $U_p$ recurrence. Thus $\bar f$ is a
nonzero simultaneous eigenvector. To lift it, take the integral
cuspidal $q$-expansion lattice in characteristic zero with the exact
character $\eta$. Its Hecke algebra generated by these operators
is finite and torsion-free over the coefficient DVR, being a
submodule of the endomorphisms of a finite free lattice. The
maximal ideal of the residual eigenvalues therefore contains a
minimal prime of characteristic zero. The corresponding characteristic-zero
eigensystem occurs on the cuspidal space: the Hecke action is faithful
and the commuting operators on its nonzero localized space have a
common eigenvector after finite scalar extension. If its first
coefficient were zero, the usual Hecke coefficient formula for
$T_n$ (including the bad-prime operators) would give $a_n=0$
for every $n$. Normalize the eigenvector at one. We obtain a cusp
eigenform $g$ satisfying
\[
       a_p(g)\equiv1+\eta(p)p^{k-1}\qquad\text{for every prime }p .
\]
Its level is exactly $C$: a form of lower new level would have
nebentype of conductor dividing that level, whereas $\eta$ has
conductor $C$. In particular $g$ is a newform. The lower bound on
$k$ is chosen once, for the fixed level, before any $k$ of the
required parity is prescribed.
% END P66E3

% BEGIN P66E4
\paragraph{P66E4. The rational Galois representation with geometric attachment.}
Take the classical representation $V_g$ in arithmetic Frobenius
convention. The Deligne--Ribet theorem gives a continuous absolutely
irreducible representation over a finite extension of $\mathbb Q_\ell$,
unramified at $p\nmid C\ell$, with trace $a_p(g)$ and determinant
$D(p)=p^{k-1}\eta(p)$ there. Put $\rho=V_g$ and write
$D=\chi_\ell^{k-1}\phi\theta$ and $\Phi=\chi_\ell^{k-1}\phi$.
Chebotarev gives $\det\rho=D$ on all group elements. Reducing a
stable lattice and applying the same density to its finite residual
image and the residual characters gives
\[
 \det(X-\bar\rho(h))=(X-1)(X-\bar\Phi(h))
 \qquad(h\in G_{\mathbb Q}).
\]
Here $\bar\theta=1$; the identity is coefficientwise and makes
no choice of residual splitting.

For the source's automorphic normalization, let $\pi_u$ be the
unitary automorphic representation of $g$. At a good prime with
arithmetic roots $\alpha,\beta$, its Satake pair is
$p^{-(k-1)/2}(\alpha,\beta)$. Use the cusp
$\Pi=\pi_u^\vee\otimes|\det|^{k/2}$. Its infinitesimal-character
entries, up to the sign convention for the determinant norm, are a
common shift by $k/2$ of $\pm(k-1)/2$, hence distinct half-integers.
At every good prime $p^{1/2}$ times its Satake pair is
$(\alpha^{-1},\beta^{-1})$, exactly the eigenvalues of
$\rho(\mathrm{Frob}_p^{-1})$ after applying $\iota$.
This proves the required regular algebraic attachment to the same
$\rho$. The classical inputs are the modular-form representation
theorem of Deligne and the characteristic-zero irreducibility
theorem of Ribet, \emph{Galois representations attached to eigenforms
with Nebentypus}, \S2.
% END P66E4

% BEGIN P66E5
\paragraph{P66E5. Ordered ordinary characters, tame inertia and the plain cusp.}
At $\ell$, P66E3 gives $a_\ell(g)\equiv1$. The Hecke polynomial
$X^2-a_\ell(g)X+\ell^{k-1}\eta(\ell)$ has one unit root
$u\equiv1$. The ordinary representation theorem for a classical
newform of tame level prime to $\ell$ gives an unramified
\emph{quotient} $\nu$ of $V_g$ with arithmetic Frobenius value $u$.
Hence on a decomposition group the upper triangular flag has
characters $(D\nu^{-1},\nu)$ in this order. Since $\phi\theta$
is unramified at $\ell$, their inertia restrictions are
$(\chi_\ell^{k-1},1)$. The reduction of $\nu$ is trivial on the
entire decomposition group, since it is unramified and its
Frobenius value reduces to one. The diagonal reductions are
therefore $(\bar\Phi,1)$ on that whole group. Conjugating the
local flag gives the assertion at each specified place above $\ell$.
We use the characteristic-zero ordinary representation theorem,
as in Skinner--Wiles, \emph{Ordinary representations and modular
forms}, opening section.

At $p\mid C$ the local conductor of the newform equals the
conductor of its central character, with positive exponent.
The two-dimensional conductor comparison in SX3 forces a principal
series with one unramified character and the other of the full
conductor, and zero monodromy. Indeed an irreducible Weil parameter
has tame codimension two, strictly greater than that of its
determinant, and on higher ramification groups determinant
codimension is at most vector codimension. A special parameter
has excess conductor as well, and two ramified characters have
sum of conductors strictly greater than the conductor of their
product. The two inducing characters have distinct inertia, so
are unlinked. Local--global compatibility for classical newforms
away from $\ell$ (Langlands--Carayol) consequently gives
\emph{actual} diagonal inertia on $V_g$, with characters
$(1,\phi\theta)$ by the determinant. The Weil--Deligne
monodromy vanishes and Frobenius semisimplification preserves
inertia. Away from $C\ell$ both characters are trivial on
inertia. The trivial valuation has trivial inertia as well.
This gives every away inertia condition. The cusp $\Pi$ of P66E4
is spherical outside $C$ (including at $\ell$) and its two
characters at the other places remain unlinked under duality
and common twist. It thus supplies the very same $\mathrm{PR}$
witness as its automorphic attachment. Together P66E1--P66E5
prove $\mathrm{PropEGP}$ over $\mathbb Q$ for the full range
of its data. This proof uses the rational newform theorems
directly at the fixed $\iota$.
% END P66E5
